\section{Decision log (abridged)}
\label{app:decisions}

The project kept a numbered decision register. The entries that bear on claims in
this paper are reproduced here; the full register, with dates and rationale for
all $23$ entries, is in \texttt{DECISIONS.md}.

\begin{table*}[t]
\centering
\small
\caption{Abridged decision register. Entries omitted here concern file layout and
tooling rather than claims.}
\label{tab:decisions}
\setlength{\tabcolsep}{4pt}
\begin{tabular}{@{}p{0.055\linewidth}p{0.86\linewidth}@{}}
\toprule
ID & Decision and reason \\
\midrule
D010 & Reject the original split. $76$ cross-split base-name groups; $63.8\%$
of evaluation images with a pHash-$\le2$ training neighbour; augmentation before
splitting confirmed in the upstream scripts. \\
D011 & Image-level and FP/image metrics blocked until the test split contains
negatives. It contained zero. \\
D012 & \texttt{May}/\texttt{Jan} files are colour RGB with marker pins, not
thermal, despite the \texttt{lwir} filename. Keep and disclose (D016). \\
D014 & Operational leak definition: exact MD5, or base key, or same-source
pHash$\le2$ \emph{and} $32\times32$ RGB MSE $\le 900$. Manual review of all $146$
candidates sets the MSE cut (same-scene $\le 817$, texture $\ge 958.6$). \\
D015 & Freeze \texttt{landmine\_combined\_v2} at 3{,}907/181/189 with $30$ eval
leaks removed; manifest digest \texttt{748322d6\ldots c80de}; true leak rate
$0/370$. \\
D017 & Disclosed protocol deviation: the upstream training script auto-evaluates
test at confidence $0.25$. No selection uses those numbers; the frozen-protocol
evaluation is a separate single pass. \\
D018 & Accept early stopping at epoch $67/75$ (best $47$) for the pilot run;
retrain at batch $8$ for comparability (D019). The pilot is not a paper result. \\
D020 & Frozen-feature MNv3 is degenerate on negatives at every resolution and
batch size ($\kappa \approx 0$); only fine-tuning the last layers separates them.
Stage A rows retained as the mandated ablation. \\
D021 & Threshold rule: sweep protection $P$, choose the largest $\tau$ within
$0.5$ points of stage-1 validation recall, then minimise validation FP/image.
All seeds select $P=0.4$. \\
D022 & Grad-CAM flags images by energy in the union of all ground-truth boxes,
not per detection; per-detection rows were an instance-count artefact on
multi-mine frames. \\
D023 & Remove two objectively wrong test boxes (tree trunk, fence); report
before/after; do not retrain (no training-label errors found). \\
\bottomrule
\end{tabular}
\end{table*}

\section{Test-touch accounting}
\label{app:touches}

Eighteen numbered test touches were logged during the original run; a further
seven were logged during the extended analysis. Grouped by purpose:

\begin{table*}[t]
\centering
\small
\caption{Test-touch accounting, grouped by purpose.}
\label{tab:touches}
\begin{tabular}{@{}p{0.10\linewidth}p{0.34\linewidth}p{0.48\linewidth}@{}}
\toprule
touches & purpose & sanctioning rule \\
\midrule
\#1--\#2 & audit read-only passes & leakage audit reads labels only \\
\#3 & MNv3 baseline, seed 0 & Phase 4 protocol eval, val-chosen thresholds \\
\#4 & YOLO seed-0 frozen-protocol eval & the single protocol-governing eval \\
\#5 & MNv3 seeds 1--2 & Phase 4 protocol eval \\
\#6--\#8 & two-stage seeds 0--2 & Phase 5 protocol eval, val-frozen \\
\#9--\#12 & Grad-CAM E1--E6 & Phase 6, read-only attribution \\
\#13--\#14 & YOLO seeds 1--2 & Phase 3 per-seed protocol eval \\
\#15--\#16 & leaky baseline, v1 split and clean subset & extended analysis \\
\#17 & YOLOv8n size point & size ablation, val-chosen threshold \\
\#18 & corrected detector-aligned two-stage & re-score, no selection \\
\#19 & YOLOv8n seeds 1--2 and YOLOv8s seeds 3--4 & seed extension, val-chosen thresholds \\
\#20 & operating-point sweep, five seeds per size & re-read of committed detection records, no selection \\
\#21 & stratified Grad-CAM re-cut & read-only attribution, strata fixed in advance \\
\#22 & threshold-selection bootstrap & re-read; validates the val-only rule, chooses nothing \\
\#23 & five-seed stage-2 audit & pipeline rerun with each seed's own detector \\
\#24 & seed-3/4 two-stage pipeline & extended analysis, val-chosen configuration \\
\#25 & embedding near-duplicate check & read-only audit, no detector evaluated \\
\bottomrule
\end{tabular}
\end{table*}

Only touch \#4 could contaminate a claim, and it is the evaluation the paper's
detector numbers come from. Touch \#18 is a correction: the original seed-1/2
two-stage runs had reused the seed-0 detector, which we detected from
byte-identical arm-A/arm-B records and corrected by re-running with each seed's
own checkpoint. No threshold was re-chosen on test; the per-seed validation
selection was rerun for the corrected detector records. Touches \#19--\#24 repeat
the same protocol on additional seeds; none of them altered a rule, a threshold
grid, or a selection criterion after looking at test.

\section{Failure gallery and reproduction commands}
\label{app:repro}

The false-positive contact sheet is reproduced in Figure~\ref{fig:failures}. The
full false-negative gallery, the audit montages, the per-image predictions, and
the Grad-CAM heatmaps are in the repository under \texttt{paper\_results/}.

\begin{figure*}[t]
\begin{verbatim}
# tables T1-T6 and figures (no GPU, no dataset required)
python src/reporting/make_tables.py

# frozen-protocol evaluation of a checkpoint
python src/evaluation/phase3_eval.py --run <run_dir> --out <out_dir>

# two-stage pipeline with each detector seed
python src/pipeline/phase5_pipeline.py --yolo <seed_best.pt> \
    --cls <stage2_best.pt> --arm-a-conf <val_conf> --out <out_dir>

# statistical analysis, bootstrap and policy rules
python src/analysis/analysis_stats.py
\end{verbatim}
\end{figure*}

Environment: Python 3.11.9, Ultralytics 8.4.117, PyTorch 2.6.0+cu124,
torchvision 0.21.0+cu124, NVIDIA RTX~4070 Laptop (8\,GB). Evaluation at
$640\times640$, NMS IoU $0.7$, \texttt{max\_det} $300$, FP32, no TTA.
